%!Mode:: "TeX:UTF-8"
%!TEX TS-program = xelatex %!Mode:: "TeX:UTF-8"

\documentclass[a4paper]{rukh-cv-zh}

\begin{document}
%
%----------------------------------------------------------------------------------------
%	LEFT COLUMN
%----------------------------------------------------------------------------------------
%
\begin{minipage}[t]{0.19\textwidth} % The left column takes up 33% of the text width of the page

\namesection{丁}{鹏}

\infosection{18611105508 \\
\href{mailto:i@dingpeng.me}{i@dingpeng.me}  \\
北京市海淀区清华大学 \\
信息大楼1-309}{\includegraphics[width=.9em]{icons/linkedin.pdf} \href{https://www.linkedin.com/in/dingpeng/zh-cn}{dingpeng} \\
\includegraphics[width=.9em]{icons/github.pdf} \href{https://github.com/rxkh}{rxkh}}

\end{minipage} % The end of the left column
\hfill
%
%----------------------------------------------------------------------------------------
%	RIGHT COLUMN
%----------------------------------------------------------------------------------------
%
\begin{minipage}[t]{0.74\textwidth} % The right column takes up 66% of the text width of the page

%------------------------------------------------
% Education
%------------------------------------------------

\section{教育}

\runsubsection{清华大学}
\descript{| 工学硕士在读，软件工程专业}
\location{2013/8 - 2016/7 | 北京}
\vspace{\topsep}
\begin{tightitemize}
\item 从事基于话题模型的事件跟踪及趋势预测研究
\end{tightitemize}

\sectionspace

%------------------------------------------------

\runsubsection{清华大学}
\descript{| 工学学士，计算机软件专业}
\location{2009/8 - 2013/7 | 北京}
\begin{tightitemize}
\item 清华大学优秀学生干部、清华大学社会工作单项优秀奖学金
\end{tightitemize}

\sectionspace

%------------------------------------------------
% Projects
%------------------------------------------------

\section{项目}

\runsubsection{CodingYoung}
\descript{| 公益编程课程}
\location{2014/4 – 2015/6 | 北京}
\begin{tightitemize}
\item 设计并开发中学生编程课程《Python编程实践》12讲
\item 在清华大学附属中学开设选修课，已帮助25名初中生学习编程思想和方法
\end{tightitemize}

\sectionspace

%------------------------------------------------

\runsubsection{ANSYS中国市场办公解决方案}
\descript{| 企业应用}
\location{2014/4 - 2015/1 | 北京}
\begin{tightitemize}
\item 开发微信服务号 (微信号: ANSYS-China)，订阅数在6个月中增长6倍
\item 开发网站监控工具，每周爬取网站内容并做出内容更新报告
\item 开发 JavaScript Web Form 工具，实现低权限下数据安全跨域传输机制
\end{tightitemize}

\sectionspace

%------------------------------------------------

\runsubsection{学生清华}
\descript{| 校园网站}
\location{2013/8 - 2014/7 | 北京}
\begin{tightitemize}
\item 学生第二课堂管理网站，独立用户数20000，平均单日访问1700
\item 主持学生活动和资源管理、新版首页等新上线功能的产品设计和技术架构
\item 负责网站开发的进度控制和团队的日常管理
\end{tightitemize}

\sectionspace

%------------------------------------------------

\runsubsection{MICAPS}
\descript{| 应用算法}
\location{2013/1 - 2013/2 | 北京}
\begin{tightitemize}
\item 天气实时监测和预报制作系统，部署于国家气象中心
\item 设计基于等值线拖拽的交互式气象预报数据编辑算法，并重写对应模块
\item 将单次编辑耗时由约10秒减少至0.1秒
\end{tightitemize}

\sectionspace

%------------------------------------------------

\runsubsection{晒课厅}
\descript{| 人人网应用}
\location{2011/5 - 2013/1 | 北京}
\begin{tightitemize}
\item 帮助清华学生在人人网上评价课程和分享课程经验
\item 上线两周5000用户，累计用户15000，人人网教育类应用排名第二达两周
\item 负责全部本科生课程数据的爬取和部分页面的前、后端实现
\end{tightitemize}

\sectionspace

%------------------------------------------------
% Experience
%------------------------------------------------

\section{经历}

\runsubsection{北京白驹谷科技有限责任公司}
\descript{| 技术负责人}
\location{2014/4 - 2015/4 | 北京}
\begin{tightitemize}
\item 网站、移动应用、微信轻应用的前、后端开发
\end{tightitemize}

\sectionspace

%------------------------------------------------

\runsubsection{清华大学成才信息中心}
\descript{| 技术负责人、辅导员}
\location{2013/8 - 2014/7 | 北京}
\begin{tightitemize}
\item 带领技术团队建设和维护新版学生清华网站
\end{tightitemize}

\sectionspace

%------------------------------------------------
% Skills
%------------------------------------------------

%\section{Skills}

%Python, Django, C/C++, C\#, Java, SQL, MongoDB

\sectionspace

\end{minipage}

\end{document}
